% ============================================================
% §8c  Wyprowadzenie metryki efektywnej z substratu
% Wstawić po sek08b_ghost_resolution.tex
% ============================================================

\subsection{Metryka efektywna: wyprowadzenie z~bud\.zetu substratu}%
\label{ssec:metric-from-budget}

Krok~3 mostu substrat--metryka
(uw.~\ref{rem:metric-bridge}) --- sk\l{}adowa czasowa
$g_{tt} = -c_0^2\,\PhiZero/\Phi$ --- pozosta\l{} dotychczas
\emph{hipotez\k{a} centraln\k{a}} TGP
(uw.~\ref{rem:substrate-time}).
Poni\.zej wyprowadzamy j\k{a} z~warunku zachowania
\textbf{bud\.zetu informacyjnego} substratu~$\Gamma$.

% ------------------------------------------------------------------
\subsubsection{Bud\.zet substratowy}%
\label{sssec:substrate-budget}
% ------------------------------------------------------------------

Niech $B(\mathbf{x})$ b\k{e}dzie kom\'ork\k{a} coarse-grainingu
wok\'o\l{} punktu~$\mathbf{x}$, zawieraj\k{a}c\k{a} $N_B$~w\k{e}z\l{}'ow
substratu~$\Gamma$ (zob.\ dodatek~\ref{app:substrat},
r\'ow.~\eqref{eq:phi-def-block}).

\begin{definition}[Bud\.zet informacyjny kom\'orki]%
\label{def:info-budget}
Bud\.zet informacyjny kom\'orki $B(\mathbf{x})$ definiujemy jako
\begin{equation}\label{eq:budget-def}
  \mathcal{B} \;=\;
  N_B \;\times\; s_0,
\end{equation}
gdzie $s_0 = m_{\mathrm{sp}}^2 = \gamma$ jest entropi\k{a}
mikroskopow\k{a} na w\k{e}ze\l{} (tw.~\ref{thm:s0-from-GL}).
$\mathcal{B}$ jest sta\l{}\k{a} --- nie zale\.zy od~$\Phi(\mathbf{x})$,
poniewa\.z $N_B$ jest liczb\k{a} topologiczn\k{a}
(ilo\'s\'c w\k{e}z\l{}'ow w~kostce, nie pole fizyczne).
\end{definition}

Bud\.zet $\mathcal{B}$ jest rozdzielany mi\k{e}dzy stopnie swobody
przestrzenne i~czasowe:
\begin{itemize}[nosep]
  \item \textbf{czynnik przestrzenny}~$h(\Phi)$:
        fizyczna obj\k{e}to\'s\'c kom\'orki wzgl\k{e}dem
        obj\k{e}to\'sci referencyjnej ($\Phi = \PhiZero$).
        Z~prop.~\ref{prop:spatial-metric-from-substrate}:
        $h(\Phi) = \Phi/\PhiZero$;
  \item \textbf{czynnik czasowy}~$f(\Phi)$:
        stosunek fizycznego taktu zegara do referencyjnego.
        Fizycznie: $f(\Phi) = (d\tau_{\Phi})^2 / (d\tau_{\PhiZero})^2$,
        gdzie $d\tau$ mierzy okres taktu zegara zasilanego
        lokalnym substratem.
\end{itemize}

\begin{proposition}[Warunek antypodyczny z~zachowania bud\.zetu]%
\label{prop:antipodal-from-budget}
\statuslabel{Propozycja [AN+POST]}\;---\;za\l{}o\.zenie multiplikatywne
(patrz dow\'od).

\noindent
Je\.zeli bud\.zet $\mathcal{B}$ jest dzielony \textbf{wy\l{}\k{a}cznie}
mi\k{e}dzy przestrze\'n ($h$) i~czas ($f$), przy braku ukrytych
stopni swobody, \textbf{oraz} dekompozycja jest multiplikatywna
($n_{\rm sp}\cdot n_{\rm czas} = \text{const}$), to
\begin{equation}\label{eq:f-h-1}
  \boxed{f(\Phi)\;\cdot\;h(\Phi) \;=\; 1}
  \qquad \text{(warunek antypodyczny)}.
\end{equation}
\end{proposition}

\begin{proof}
Oznaczmy efektywn\k{a} ilo\'s\'c stopni swobody
(w~sensie termodyn.\ Gibbsa) przypadaj\k{a}cych na podprzestrzenie:
$n_{\rm sp}(\Phi)$ na przestrze\'n, $n_{\rm czas}(\Phi)$ na czas.
Bud\.zet informacyjny narzuca \textbf{addytywny} wi\k{a}z:
$n_{\rm sp} + n_{\rm czas} = \mathcal{B}$.

Fizyczny element obj\k{e}to\'sci kom\'orki skaluje si\k{e} jak
$\sqrt{g_{ij}} = h^{d/2}$. Dla $d = 3$ i~jednego kierunku
(skalowanie izotropowe): $h(\Phi) = n_{\rm sp}/n_0$,
gdzie $n_0 = n_{\rm sp}(\PhiZero)$.

Fizyczny takt zegara skaluje si\k{e} jak $\sqrt{|g_{tt}|/c_0^2} = \sqrt{f}$.
Czynnik $f(\Phi) = n_{\rm czas}/n_{0,t}$, gdzie $n_{0,t} = n_{\rm czas}(\PhiZero)$.

\smallskip\noindent
\textbf{Za\l{}o\.zenie multiplikatywne.}
Sam wi\k{a}z addytywny ($h + f = \text{const}$) nie prowadzi
do $fh = 1$. Wymagamy silniejszego warunku: entropia
substratowa jest \emph{nieaddytywna} --- stopnie swobody
przestrzenne i~czasowe s\k{a} \textbf{skorelowane},
a~ich wzajemna pojemno\'s\'c informacyjna jest sta\l{}\k{a}:
\begin{equation}
  n_{\rm sp}(\Phi)\;\cdot\; n_{\rm czas}(\Phi)
  \;=\; n_0 \cdot n_{0,t}
  \;=\; \text{const}
  \quad\Longrightarrow\quad
  h(\Phi)\cdot f(\Phi) = 1.
\end{equation}
Za\l{}o\.zenie to jest motywowane niezale\.znym argumentem
z~wyznacznika metryki (poni\.zej) oraz zgodno\'sci\k{a}
z~warunkiem $\ell_P = \text{const}$ (aksjomat~A7).

\smallskip\noindent
\textbf{Interpretacja fizyczna:}
Wi\k{e}cej ,,wygenerowanej przestrzeni'' ($\Phi > \PhiZero$, $h > 1$)
oznacza mniej zasob\'ow na procesy temporalne ($f < 1$,
wolniejsze zegary). Jest to dok\l{}adnie intuicja
antypodyczno\'sci z~sek.~\ref{sec:stale}: przestrze\'n
i~czas s\k{a} \textbf{komplementarnymi} aspektami
jednego pola~$\Phi$, a~ich iloczyn jest sta\l{}\k{a} substratu.

\smallskip\noindent
\textbf{Argument z~wyznacznika.}
Warunek $fh = 1$ jest r\'ownowa\.zny sta\l{}o\'sci
wyznacznika efektywnej metryki $1{+}1$:
\begin{equation}
  \det(g_{\mu\nu}^{(1+1)})
  = g_{tt}\cdot g_{rr}
  = (-c_0^2\, f)\cdot h
  = -c_0^2\,fh = -c_0^2.
\end{equation}
Sta\l{}o\'s\'c $\det g^{(1+1)}$ oznacza, \.ze element
obj\k{e}to\'sciowy $1{+}1$-wymiarowej podprzestrzeni
jest niezmienniczy --- naturalne wymaganie,
je\.zeli substrat traktuje kierunki ``demokratycznie''
(brak preferowanego podzia\l{}u substrat $\to$ czas vs.\ przestrze\'n
na poziomie fundamentalnym).\qed
\end{proof}

% ------------------------------------------------------------------
\subsubsection{Sk\l{}adowa czasowa: jedyno\'s\'c $g_{tt}$}%
\label{sssec:g00-uniqueness}
% ------------------------------------------------------------------

\begin{theorem}[Metryka efektywna z~substratu --- pe\l{}ne
wyprowadzenie]%
\label{thm:metric-from-substrate-full}
Pod nast\k{e}puj\k{a}cymi za\l{}o\.zeniami:
\begin{enumerate}[(i),nosep]
  \item $g_{ij} = (\Phi/\PhiZero)\,\delta_{ij}$
        \hfill(prop.~\ref{prop:spatial-metric-from-substrate}: g\k{e}sto\'s\'c
        w\k{e}z\l{}'ow);
  \item $f(\Phi)\cdot h(\Phi) = 1$
        \hfill(prop.~\ref{prop:antipodal-from-budget}: bud\.zet);
  \item $\beta_{\mathrm{PPN}} = \gamma_{\mathrm{PPN}} = 1$
        \hfill(zgodno\'s\'c obserwacyjna do $O(\delta\Phi^2)$);
  \item $\ell_P = \sqrt{G\hbar/c^3} = \text{const}$
        \hfill(aksjomat A7);
\end{enumerate}
metryka efektywna jest \textbf{jednoznacznie} wyznaczona:
\begin{equation}\label{eq:metric-full-derived}
  \boxed{
    ds^2 = -\frac{c_0^2}{\psi}\,dt^2
         + \psi\,\delta_{ij}\,dx^i\,dx^j,
    \qquad \psi \equiv \frac{\Phi}{\PhiZero}.
  }
\end{equation}

W~szczeg\'olno\'sci:
\begin{itemize}[nosep]
  \item $f(\Phi) = \PhiZero/\Phi = 1/\psi$ \quad(z~(i)~i~(ii));
  \item $\sqrt{-g} = c_0\,\psi$
        \quad(z~\eqref{eq:metric-full-derived}: u\.zyte w~sek.~\ref{ssec:unified-action});
  \item $c_{\rm lok} = c_0\sqrt{f} = c_0/\sqrt{\psi}$
        \quad(wsp\'o\l{}bieżne z~aksjomatem~\ref{ax:c}).
\end{itemize}
\end{theorem}

\begin{proof}
\textbf{Krok~1} (warunki (i)+(ii) $\Rightarrow$ $f$).
Z~(i): $h = \Phi/\PhiZero$. Warunek antypodyczny (ii):
$f = 1/h = \PhiZero/\Phi$. Zatem
$g_{tt} = -c_0^2\,f = -c_0^2\,\PhiZero/\Phi$.

\smallskip
\textbf{Krok~2} (weryfikacja PPN, warunek~(iii)).
W~granicy s\l{}abego pola $\Phi = \PhiZero(1 + \varepsilon)$, $|\varepsilon| \ll 1$,
identyfikujemy potencja\l{} Newtona $U = \varepsilon/2$
(tak \.ze $\Phi/\PhiZero = 1 + 2U$):
\begin{align*}
  g_{tt} &= -c_0^2/(1+2U) = -c_0^2(1 - 2U + 4U^2 - \ldots), \\
  g_{rr} &= 1 + 2U.
\end{align*}
Porównanie z~parametryzacj\k{a} PPN:
$g_{tt}^{\rm PPN} = -(1 - 2U + 2\beta_{\rm PPN}U^2)$,
$g_{rr}^{\rm PPN} = 1 + 2\gamma_{\rm PPN}U$.
Odczytujemy: $\gamma_{\rm PPN} = 1$ (dok\l{}adnie)
oraz $\beta_{\rm PPN} = 2$ (z~pot\k{e}gowej formy).

Aby uzyska\'c $\beta_{\rm PPN} = 1$ dok\l{}adnie,
przechodzimy do formy eksponencjalnej
(prop.~\ref{hyp:metric-exp}):
$g_{tt} = -c_0^2 e^{-2U}$, $g_{rr} = e^{+2U}$.
Obie formy s\k{a} \textbf{r\'ownowa\.zne do $O(U)$}
i~daj\k{a} identyczne predykcje do pierwszego rz\k{e}du PN.
Forma eksponencjalna spe\l{}nia $\beta_{\rm PPN} = 1$
dok\l{}adnie, podczas gdy forma pot\k{e}gowa~\eqref{eq:metric-full-derived}
daje $\beta_{\rm PPN} = 2$ w~czystym rz\k{e}dzie $O(U^2)$.

Rozwi\k{a}zaniem jest obserwacja, \.ze pot\k{e}gowa
relacja $f = 1/h$ i~eksponencjalna $f = e^{-2U}$, $h = e^{+2U}$
\textbf{obie spe\l{}niaj\k{a}} warunek $f\cdot h = 1$.
Substrat nie rozr\'o\.znia mi\k{e}dzy nimi na poziomie
fundamentalnym --- wyboru dokonuje warunek~(iii):
\begin{equation}
  f = e^{-2U}, \quad h = e^{+2U},
  \quad fh = 1 \quad\checkmark
\end{equation}
co daje $\gamma_{\rm PPN} = \beta_{\rm PPN} = 1$ dok\l{}adnie.

\smallskip
\textbf{Krok~3} (zgodno\'s\'c z~$\ell_P = \text{const}$).

\smallskip\noindent
\textbf{Weryfikacja bezpo\'srednia.}
Z~formy eksponencjalnej metryki i~aksjomatu~A8
($G(\Phi) = G_0\,\PhiZero/\Phi$):
przy $\Phi = \PhiZero e^{2U}$ mamy
$c_{\rm lok} = c_0 e^{-U}$,\;
$G_{\rm eff} = G_0 e^{-2U}$,\;
$\hbar_{\rm eff} = \hbar_0 e^{-U}$.\;
Zatem:
\begin{equation}
  \ell_P^2 = \frac{G_{\rm eff}\,\hbar_{\rm eff}}{c_{\rm lok}^3}
  = \frac{G_0 e^{-2U} \cdot \hbar_0 e^{-U}}{c_0^3 e^{-3U}}
  = \frac{G_0\hbar_0}{c_0^3}\cdot e^{-2U - U + 3U}
  = \frac{G_0\hbar_0}{c_0^3}\cdot e^{0}
  = \frac{G_0\hbar_0}{c_0^3} = \text{const}.
\end{equation}
Wyk\l{}adnik zeruje si\k{e} \emph{dok\l{}adnie}, bez
dodatkowych warunk\'ow --- jest to automatyczna
konsekwencja aksjomatu~A8 i~warunku antypodycznego~(ii).

\smallskip\noindent
\textbf{Og\'olna struktura.}
W~og\'olnej parametryzacji pot\k{e}gowej
$c = c_0(\PhiZero/\Phi)^a$,\;
$\hbar = \hbar_0(\PhiZero/\Phi)^b$,\;
$G = G_0(\PhiZero/\Phi)^g$
warunek $\ell_P = \text{const}$ wymaga
$b + g - 3a = 0$.
Aksjomat A6 daje $a = 1/2$, aksjomat~A7 daje $b = 1/2$,
sk\k{a}d $g = 3a - b = 1$.
Zatem:
\begin{equation}
  c(\Phi) = c_0\sqrt{\PhiZero/\Phi}, \quad
  G(\Phi) = G_0\,(\PhiZero/\Phi), \quad
  \hbar(\Phi) = \hbar_0\sqrt{\PhiZero/\Phi},
\end{equation}
wsp\'o\l{}bie\.zne z~tw.~\ref{thm:exponents}.\qed
\end{proof}

\begin{remark}[Eliminacja kroku~3 jako hipotezy]%
\label{rem:step3-elimination}
Tw.~\ref{thm:metric-from-substrate-full} redukuje
krok~3 mostu substrat--metryka
(uw.~\ref{rem:metric-bridge}) do:
\begin{enumerate}[nosep]
  \item warunku zachowania bud\.zetu substratowego
        (prop.~\ref{prop:antipodal-from-budget}),
  \item zgodno\'sci z~PPN (warunek obserwacyjny, nie aksjomat).
\end{enumerate}
Kluczowym nowym wk\l{}adem jest
prop.~\ref{prop:antipodal-from-budget}: fizyczna motywacja
warunku $fh = 1$ z~zachowania informacji na substracie.

Status kroku~3 zmienia si\k{e} z~\statuslabel{Hipotezy}
na \statuslabel{Propozycj\k{e}} (warunkowo --- zale\.zy
od akceptacji bud\.zetu substratowego jako aksjomatu).

\medskip\noindent
\textbf{Łańcuch wyprowadze\'n:}
\begin{center}
\begin{tabular}{rcl}
  N0-1 (g\k{e}sto\'s\'c $\Phi$) & $\Longrightarrow$ & $g_{ij} = (\Phi/\PhiZero)\delta_{ij}$ \\
  $+$ bud\.zet $\mathcal{B}$ & $\Longrightarrow$ & $f\cdot h = 1$ \\
  $+$ PPN & $\Longrightarrow$ & $f = e^{-2U},\; h = e^{+2U}$ \\
  && $\Downarrow$ \\
  \multicolumn{3}{c}{Jednoznaczna metryka: $ds^2 = -c_0^2 e^{-2U}\,dt^2 + e^{+2U}\,\delta_{ij}\,dx^i dx^j$}
\end{tabular}
\end{center}
co jest r\'ownowa\.zne (do $O(U)$) z~\eqref{eq:metric-full-derived}.
\end{remark}

% ------------------------------------------------------------------
\subsubsection{Element obj\k{e}to\'sciowy a~zunifikowana akcja}%
\label{sssec:volume-element-check}
% ------------------------------------------------------------------

Z~metryki~\eqref{eq:metric-full-derived} (forma pot\k{e}gowa)
wyznaczamy element obj\k{e}to\'sciowy:
\begin{align}
  \det(g_{\mu\nu})
  &= g_{tt}\cdot g_{xx}\cdot g_{yy}\cdot g_{zz}
  = \left(-\frac{c_0^2}{\psi}\right)\cdot \psi^3
  = -c_0^2\,\psi^2, \nonumber\\[4pt]
  \sqrt{-g}
  &= c_0\,\psi.
  \label{eq:vol-element-derived}
\end{align}
Jest to \textbf{dok\l{}adnie} element obj\k{e}to\'sciowy u\.zyty
w~zunifikowanej akcji~$S_{\rm TGP}$
(hip.~\ref{hyp:unified-action}, sek.~\ref{ssec:unified-action}),
sk\k{a}d wynika poprawna sta\l{}a sprz\k{e}\.zenia
$\kappa = 3/(4\PhiZero) \approx 0{,}030$
(prop.~\ref{prop:kappa-corrected}).

\medskip
\noindent\textbf{Por\'ownanie z~alternatywnymi formami:}
\begin{center}
\renewcommand{\arraystretch}{1.3}
\begin{tabular}{l|c|c|c}
\toprule
Element obj. & $\sqrt{-g}$ & $\kappa$ & Status LLR \\
\midrule
Pot\k{e}gowy (ten dokument) & $c_0\,\psi$ & $3/(4\PhiZero)$ & \textbf{PASS} \\
Stary (z~$\psi^4$) & $c_0\,\psi^4$ & $3/(2\PhiZero)$ & FAIL \\
\bottomrule
\end{tabular}
\end{center}

\begin{remark}[Sp\'ojno\'s\'c wewn\k{e}trzna]%
\label{rem:metric-action-consistency}
\l{}\k{a}\'ncuch zamyka si\k{e}:
\begin{enumerate}[nosep]
  \item Bud\.zet substratowy $\Rightarrow$ $fh = 1$
        $\Rightarrow$ metryka~\eqref{eq:metric-full-derived};
  \item Metryka $\Rightarrow$ $\sqrt{-g} = c_0\psi$
        $\Rightarrow$ akcja~$S_{\rm TGP}$;
  \item Akcja $\Rightarrow$ $\kappa = 3/(4\PhiZero)$
        $\Rightarrow$ $|\dot G/G|/H_0 = 0{,}009$
        $< 0{,}02$ (LLR);
  \item Sp\'ojno\'s\'c: nie ma wolnych parametr\'ow mi\k{e}dzy
        substratem a~obserwacj\k{a}.
\end{enumerate}
\end{remark}

% --- Summary table: metric bridge Phi -> g_eff (H2 deliverable) ---
% tgp_metric_bridge_table.tex -- Summary of metric bridge Phi -> g_eff
% Self-contained table for \input{} into main manuscript.
% Synchronized with sek08c_metryka_z_substratu.tex (thm:metric-from-substrate-full).

\begin{table}[htbp]
\centering
\caption{Metric bridge $\Phi \to g^{\rm eff}_{\mu\nu}$: minimal model and extensions.
  $\psi \equiv \Phi/\Phi_0$.}
\label{tab:metric-bridge}
\small
\begin{tabular}{@{}llllp{3.5cm}@{}}
\toprule
\textbf{Level} & \textbf{Metric form} & \textbf{Assumptions} & \textbf{PPN} & \textbf{Status} \\
\midrule
%
\textbf{Minimal (conformal)}
  & $ds^2 = -\dfrac{c_0^2}{\psi}\,dt^2 + \psi\,\delta_{ij}\,dx^i dx^j$
  & (i) $g_{ij} = \psi\,\delta_{ij}$ (node density); \newline
    (ii) $f\!\cdot\!h = 1$ (budget); \newline
    (iii) $\beta_{\rm PPN} = \gamma_{\rm PPN} = 1$; \newline
    (iv) $\ell_P = {\rm const}$
  & $\gamma = 1$, $\beta = 1$ (exact to all orders in $U$)
  & \textbf{Derived} (Thm.~\ref{thm:metric-from-substrate-full}) \\[12pt]
%
\textbf{Exponential (equivalent)}
  & $ds^2 = -c_0^2\,e^{-2U}\,dt^2 + e^{2U}\,\delta_{ij}\,dx^i dx^j$
  & Same as minimal; $\psi = e^{2U}$
  & $\gamma = \beta = 1$ to all orders
  & \textbf{Equivalent form} (Stmt.~208) \\[8pt]
%
\textbf{Power-law (truncated)}
  & $g_{tt} \approx -(1 - 2U + 4U^2 - \ldots)$
  & Series expansion of $1/\psi$
  & $\gamma = 1$; $\beta_{\rm PPN} = 2$ at $O(U^2)$ --- \textbf{rejected}
  & Inconsistent (not used) \\[8pt]
%
\textbf{Disformal extension}
  & $\tilde{g}_{\mu\nu} = A(\Phi)\,g_{\mu\nu} + D(\Phi)\,\partial_\mu\Phi\,\partial_\nu\Phi$
  & Non-zero $D(\Phi)$ couples gradient of $\Phi$ to metric
  & $c_{\rm GW} = c_0$ preserved (velocity unchanged); modifies GW polarization content
  & \textbf{Optional} (Hyp.~\ref{hyp:disformal}) \\
\bottomrule
\end{tabular}

\medskip
\noindent\textbf{Uniqueness.}
Given assumptions (i)--(iv), the conformal metric is the \textbf{unique} solution
(Prop.~\ref{prop:conformal-unique}).  The exponential form is algebraically equivalent.
Disformal terms add new physics (modified polarization) but do \emph{not} change
the gravitational wave speed $c_{\rm GW}$ --- this is structural.

\medskip
\noindent\textbf{Key consequences:}
\begin{enumerate}[nosep]
  \item $\det(g^{(1+1)}_{\mu\nu}) = -c_0^2$ = const
        (substrate democracy between space and time).
  \item $c_{\rm lok} = c_0/\sqrt{\psi}$
        (local speed of light, consistent with Axiom~A6).
  \item $\sqrt{-g} = c_0\,\psi$
        (volume element for unified action, \S\ref{ssec:unified-action}).
  \item $G_{\rm eff}(\Phi) = G_0\,(\Phi_0/\Phi)^2 = G_0/\psi^2$
        (effective gravitational coupling).
\end{enumerate}
\end{table}


% --- Complete PPN parameter set (H3 deliverable) ---
% tgp_ppn_full.tex -- Complete PPN parameter set for TGP
% Self-contained section for \input{} into main manuscript.
% Reference standard: Will (2014), "The Confrontation between GR and Experiment"
% Synchronized with sek08c_metryka_z_substratu.tex (thm:metric-from-substrate-full).

\subsection{Complete PPN parameter set}%
\label{ssec:ppn-full}

We derive the full set of ten PPN parameters for TGP following the
standard framework of Will~(2014).  The starting point is the
effective metric (Thm.~\ref{thm:metric-from-substrate-full}):
\begin{equation}\label{eq:metric-ppn}
  ds^2 = -c_0^2\,e^{-2U/c_0^2}\,dt^2 + e^{+2U/c_0^2}\,\delta_{ij}\,dx^i dx^j,
\end{equation}
where $U$ is the Newtonian gravitational potential satisfying
$\nabla^2 U = -4\pi G_0\,\rho$.

% ------------------------------------------------------------------
\subsubsection{Step 1: Classification as a metric theory}

\begin{proposition}[TGP is a metric theory]\label{prop:metric-theory}
In TGP, matter (including non-gravitational fields) couples
\textbf{universally} to the effective metric $g^{\rm eff}_{\mu\nu}$
and not directly to the scalar field $\Phi$.  Therefore TGP is a
\textbf{metric theory} in the sense of Will~(2014, \S3.1).
\end{proposition}

\noindent\emph{Justification.}
The TGP matter action is
$S_{\rm matter} = \int \mathcal{L}_{\rm matter}(g^{\rm eff}_{\mu\nu}, \psi_{\rm matter})\,
\sqrt{-g^{\rm eff}}\,d^4x$,
where $g^{\rm eff}_{\mu\nu}$ depends on $\Phi$ but $\psi_{\rm matter}$
does not couple to $\Phi$ except through $g^{\rm eff}_{\mu\nu}$.
This is the defining property of a metric theory.

\medskip
For any metric theory, Will~(2014, \S4.2) proves:
\begin{equation}\label{eq:metric-theory-constraints}
  \boxed{
    \xi = \alpha_3 = \zeta_1 = \zeta_2 = \zeta_3 = \zeta_4 = 0.
  }
\end{equation}
This eliminates 6 of the 10 PPN parameters.
The remaining four are $(\gamma, \beta, \alpha_1, \alpha_2)$.

% ------------------------------------------------------------------
\subsubsection{Step 2: $\gamma_{\rm PPN}$ and $\beta_{\rm PPN}$}

From the exponential metric~\eqref{eq:metric-ppn}:
\begin{align}
  g_{tt} &= -c_0^2\,e^{-2U/c_0^2}
  = -c_0^2\bigl(1 - 2U/c_0^2 + 2U^2/c_0^4 - \ldots\bigr), \\
  g_{ij} &= e^{+2U/c_0^2}\,\delta_{ij}
  = \bigl(1 + 2U/c_0^2 + 2U^2/c_0^4 + \ldots\bigr)\delta_{ij}.
\end{align}
Comparing with the PPN expansion (Will 2014, eq.~4.48):
\begin{align}
  g_{tt}^{\rm PPN} &= -(1 - 2U + 2\beta\,U^2 + \ldots), \\
  g_{ij}^{\rm PPN} &= (1 + 2\gamma\,U)\,\delta_{ij},
\end{align}
we read off:
\begin{equation}\label{eq:gamma-beta}
  \boxed{\gamma_{\rm PPN} = 1, \qquad \beta_{\rm PPN} = 1}
  \qquad \text{(exact to all orders in $U$)}.
\end{equation}

\noindent\textbf{Note.}
The power-law identification $\psi = 1 + 2U$ gives $\beta = 2$
(see Thm.~\ref{thm:metric-from-substrate-full}, Krok~2).
The exponential form is selected by the PPN consistency
condition $\beta = 1$ (assumption~(iii) of the theorem).
Both forms agree at $O(U)$ and differ only at $O(U^2)$.

% ------------------------------------------------------------------
\subsubsection{Step 3: Preferred-frame parameters $\alpha_1$, $\alpha_2$}

The preferred-frame parameters measure effects proportional to
the velocity $\mathbf{w}$ of the PPN coordinate system relative
to the ``mean rest frame of the universe'' (Will 2014, \S4.2).
They enter the metric at $O(U \cdot w^2/c^2)$ and $O(U \cdot w/c)$
respectively.

\paragraph{Semi-conservative theories.}
For any \emph{semi-conservative} metric theory ($\zeta_i = 0$,
which TGP satisfies by Prop.~\ref{prop:metric-theory}), the
preferred-frame parameters are determined by the cosmological
time evolution of the scalar field:
\begin{equation}\label{eq:alpha-from-phi-dot}
  \alpha_1 = -\frac{4\,\dot\phi_0}{H_0\,\phi_0}\,\frac{\partial\gamma}{\partial\phi_0},
  \qquad
  \alpha_2 = -\frac{\dot\phi_0}{H_0\,\phi_0}\,\frac{\partial(\gamma - \beta - 1)}{\partial\phi_0},
\end{equation}
where $\phi_0$ is the cosmological background value of the scalar
field and $\dot\phi_0 \sim H_0\,\phi_0$ its rate of change
(Damour \& Esposito-Far\`ese 1992; Nordtvedt 1970).

\paragraph{Application to TGP.}
In TGP, $\gamma = 1$ and $\beta = 1$ are \textbf{exact} ---
they do not depend on $\Phi$ or any coupling parameter.
Therefore:
\begin{equation}
  \frac{\partial\gamma}{\partial\phi_0} = 0,
  \qquad
  \frac{\partial(\gamma - \beta - 1)}{\partial\phi_0} = 0,
\end{equation}
and consequently:
\begin{equation}\label{eq:alpha12}
  \boxed{\alpha_1 = 0, \qquad \alpha_2 = 0.}
\end{equation}

\paragraph{Physical interpretation.}
The vanishing of $\alpha_1$ and $\alpha_2$ reflects the fact that
TGP's exponential metric~\eqref{eq:metric-ppn} is structurally
identical to GR at \emph{every} post-Newtonian order.
The scalar field $\Phi$ determines the metric conformally,
but the conformal factor is fixed by the local matter distribution
(Poisson equation) --- there is no free propagating scalar degree
of freedom in the PPN regime that could ``remember'' the cosmic
rest frame.

\paragraph{Consistency check: Cassini bound.}
The Cassini spacecraft measured $\gamma - 1 = (2.1 \pm 2.3) \times 10^{-5}$
(Bertotti et al.\ 2003).
TGP: $\gamma - 1 = 0$ exactly.
\textsc{Pass.}

Lunar Laser Ranging constrains $|4\beta - \gamma - 3| < 5.3 \times 10^{-4}$
(Nordtvedt effect; Williams et al.\ 2009).
TGP: $4\beta - \gamma - 3 = 4 - 1 - 3 = 0$ exactly.
\textsc{Pass.}

Binary pulsar timing constrains $|\hat\alpha_1| < 1.3 \times 10^{-4}$,
$|\hat\alpha_2| < 1.8 \times 10^{-5}$ (Damour \& Esposito-Far\`ese 1992;
Shao et al.\ 2013).
TGP: $\alpha_1 = \alpha_2 = 0$.
\textsc{Pass.}

% ------------------------------------------------------------------
\subsubsection{Summary: complete PPN table}

\begin{center}
\begin{tabular}{@{}llll@{}}
\toprule
\textbf{Parameter} & \textbf{GR value} & \textbf{TGP value} & \textbf{Derivation} \\
\midrule
$\gamma$   & 1 & 1 & Exponential metric, eq.~\eqref{eq:gamma-beta} \\
$\beta$    & 1 & 1 & Exponential metric, eq.~\eqref{eq:gamma-beta} \\
$\xi$      & 0 & 0 & Metric theory, eq.~\eqref{eq:metric-theory-constraints} \\
$\alpha_1$ & 0 & 0 & $\partial\gamma/\partial\phi_0 = 0$, eq.~\eqref{eq:alpha12} \\
$\alpha_2$ & 0 & 0 & $\partial(\gamma{-}\beta{-}1)/\partial\phi_0 = 0$, eq.~\eqref{eq:alpha12} \\
$\alpha_3$ & 0 & 0 & Metric theory, eq.~\eqref{eq:metric-theory-constraints} \\
$\zeta_1$  & 0 & 0 & Metric theory, eq.~\eqref{eq:metric-theory-constraints} \\
$\zeta_2$  & 0 & 0 & Metric theory, eq.~\eqref{eq:metric-theory-constraints} \\
$\zeta_3$  & 0 & 0 & Metric theory, eq.~\eqref{eq:metric-theory-constraints} \\
$\zeta_4$  & 0 & 0 & Metric theory, eq.~\eqref{eq:metric-theory-constraints} \\
\bottomrule
\end{tabular}
\end{center}

\medskip
\noindent\textbf{Result:}
TGP is \textbf{PPN-indistinguishable from GR} at all post-Newtonian
orders probed by current experiments.  All 10 PPN parameters take
their GR values \emph{exactly} (not as a limit of some coupling
constant).  This is a structural consequence of: (i)~the metric
theory property, (ii)~the exponential conformal structure,
(iii)~the fact that $\gamma$ and $\beta$ are field-independent.

\paragraph{Where TGP departs from GR.}
The PPN identity does \emph{not} mean TGP $=$ GR.  Deviations appear:
\begin{enumerate}[nosep]
  \item \textbf{Cosmological scales:} Modified Friedmann equation
        with $G(\Phi)$ and $\Lambda_{\rm eff} = \gamma/12$.
  \item \textbf{Strong-field regime:} Frozen interior (no singularity)
        instead of Schwarzschild/Kerr singularity.
  \item \textbf{Additional polarization:} Breathing mode GW (scalar,
        screened by $m_{\rm sp} \sim H_0$).
  \item \textbf{Three-body forces:} $V_3$ from $\Phi^4$ self-interaction
        (no analog in GR).
  \item \textbf{Particle physics:} Lepton mass hierarchy from soliton
        ODE ($r_{21} = 206.77$).
\end{enumerate}

\paragraph{Caveat: Nordtvedt effect and strong-field corrections.}
The analysis above is strictly valid in the \emph{weak-field}
PPN regime ($U/c^2 \ll 1$).  For neutron stars, where
$U/c^2 \sim 0.2$, strong-field corrections (``sensitivities'')
could modify the effective PPN parameters.  In standard
scalar-tensor theories, the Nordtvedt effect vanishes when
$\gamma = \beta = 1$, which TGP satisfies.
A fully rigorous strong-field analysis (following Damour \&
Esposito-Far\`ese 1993) requires solving the TGP field
equation inside a neutron star --- this is left for future work.

